\begin{abstract}
	We examine whether Low-Rank Adaptation (LoRA) can reduce the cost of
	fine-tuning while retaining useful task performance, and discuss
	orthogonal LoRA (O-LoRA) for continual learning. The report reviews
	parameter-efficient fine-tuning, derives trainable-parameter counts
	for a $1024 \times 1024$ weight matrix, and compares LoRA with full
	fine-tuning of DistilRoBERTa on four-class AG News. In the worked
	example, LoRA trains 8,192 parameters at rank 4, a 99.22\% reduction
	from full fine-tuning of that matrix. In the experiment, adapters on
	the query and value layers and the classifier account for 0.8944\% of
	the model's parameters; on 7,600 test examples, the model achieves
	93.76\% accuracy and 93.76\% macro-F1. Full fine-tuning reaches
	95.03\% on both measures. The LoRA run records lower training wall
	time and peak allocated GPU memory, with accuracy 1.26 percentage
	points below the baseline in these runs.
	LoRA is therefore a practical option when reducing trainable
	parameters matters, while model quality and total resource use still
	require task-specific evaluation.
\end{abstract}
